\section{Background and Related Work}
\label{sec:related_work}


\subsection{LLM Tool-Use and Function-Calling}

Early examples of LLM tool-use, i.e., explicitly generating calls to external tools, which are then executed, were in Retrieval-Augmented Generation \citep{lewis2020retri},
where the model is allowed to retrieve information from an external source (e.g., a database) before yielding a response.
Since then, various LLM tool-use approaches have been explored in the literature 
\citep{yao2023react,schick2023toolf,qin2023tool}. The most widespread approach---which has been implemented in a very similar fashion into the majority of open- and closed-source LLMs
\citep{carrigan2024tool,openai2025fc,gemini2025fc}---is the Function-Calling (FC) interface.
In this framework, 
the LLM is explicitly provided with a list of functions in context.
Then, the LLM can call one or more of those functions to address the user's request. 








To evaluate tool-use capabilities of LLMs, many benchmarks have been proposed 
\citep{bfcl,qin2024tooll,li2023api-b,lu2024tools}.
However, all of them focus primarily on the accuracy-oriented performance of LLMs on these tool-use-related tasks. 
While accuracy is important, it is an orthogonal question to UQ, which tells us which function calls the LLM is uncertain about. 
This, in turn, can be used to refuse to execute certain  (potentially erroneous) function calls. 
This aspect, which is crucial from the reliability perspective, to the best of our knowledge, so far has been neglected in tool-use evaluation.






\subsection{Uncertainty Quantification Methods for Natural Language Generation}
\label{sec:uq-methods}

An established way to prevent incorrect model outputs before they can cause harm is to detect them via uncertainty estimates \citep{chow1957optim}. 
In problems like classification, where the output space is confined and well-structured, this can be as simple as returning the negative log likelihood (NLL) of the predicted class %
\citep{galil2023what}. When working with language models, there are at least two additional challenges.
First, rather than a single class label with a single NLL value, we now have sequences of tokens, each with their own NLL value. Second, multiple distinct sequences can have the same semantic meaning: the probability of generating a correct response is the probability of generating any one of a large set of semantically equivalent correct responses \citep{kuhn2023semantic}.
In FC, such semantic equivalences might stem from, e.g., additional whitespaces or a different order of arguments. There are two main families of UQ methods that tackle these problems. Below we give an overview of both and refer to \Cref{app:uq} for more details.

The first family uses only a single greedily decoded sequence (which is likely close to being the highest-probability sequence), 
and aggregates its individual token uncertainties, in different ways: 
\textbf{MAX} reports the highest negative log-likelihood (NLL) of all predicted tokens (i.e., highest per token uncertainty), 
\textbf{AVG} the average,
and \textbf{G-NLL} the sum \citep{aichberger2024rethi, fadeeva2023, vashurin2025benchmarkinguncertaintyquantificationmethods}. 
In addition to these,
% uncertainty estimators, 
we also evaluate the na\"{i}ve baseline \textbf{LEN} which reports the number of tokens of the greedy answer, as correctness metrics might spuriously correlate with shorter answers \citep{santilli2024on}. 
In case our benchmarking metrics were affected, it would be reflected in a high performance of this baseline.

The second family samples multiple sequences from the LLM to estimate the entropy of the LLM's distribution over whole sequences. 
A simple way of doing this is to directly estimate the Predictive Entropy (\textbf{PE}) of the distribution, based on the individual sequence's aggregated probabilities. 
However, this also captures uncertainty arising from sequences that are semantically identical. 
To account for this, Semantic Entropy (SE) \citep{kuhn2023semantic,farquhar2024detec} clusters samples that are different in form but have the same meaning, then estimates the entropy of the distribution of semantically different clusters. 
The original implementation of SE uses LLM-based entailment methods to cluster semantically-equivalent generations.
However, in initial experiments we found this to lead to bad results for the highly structured FC outputs (see \Cref{app:SEentail}), and furthermore it is of course computationally expensive. 
Thus, we use exact string matching (EXM), where two sequences get assigned the same cluster if they are exactly equal. 
Once the clusters are determined, we compute the entropy over the cluster distribution, weighing each cluster (1) by its relative frequency among all sequences (\textbf{DSE}$_{EXM}$) or (2) by the cluster probabilities computed as the sum of probabilities of its member sequences (\textbf{SE}$_{EXM}$), both following \citet{farquhar2024detec}. 

Lastly, we also include \textbf{P}$_{true}$ \citep{kadavath2022languagemodelsmostlyknow}, a prompting-based method, see \Cref{app:uq}.









